\documentclass[14pt,a4paper,oneside]{extbook}
\usepackage[utf8]{vietnam}
\usepackage{amsmath,amssymb,mathrsfs,fancyhdr,enumerate,multirow,makecell}
\usepackage{tkz-euclide,tikz-3dplot,tikz,tkz-tab}
\usetikzlibrary{arrows.meta,calc,intersections,angles,quotes,backgrounds,shapes.geometric,patterns,through,shadows,positioning}
\usepackage{yhmath}
\usepackage[top=0.7cm, bottom=0.7cm, left=1.0cm, right=1.0cm, headheight=16pt, headsep=3mm, footskip=5mm]{geometry}
\usepackage[hidelinks,unicode]{hyperref}
\usepackage{qrcode}
\usepackage[many]{tcolorbox}
\usepackage{fontawesome}

\pagestyle{fancy}
\fancyhf{}
\fancyhead[L]{\footnotesize\bfseries\color{teal!85!black}\faBook~SỔ TAY KIẾN THỨC TOÁN 5}
\fancyhead[R]{\footnotesize\bfseries\color{blue!85!black}Lớp Toán Thầy Lĩnh -- loptoanthaylinh.vn}
\fancyfoot[C]{\footnotesize\bfseries\color{gray!80!black}-- Sổ tay kiến thức: Bài 71. Ôn tập hình học --}
\renewcommand{\headrulewidth}{0.4pt}

\begin{document}

% --- BANNER TIÊU ĐỀ BÀI HỌC ---
\noindent
\begin{tcolorbox}[
  enhanced,
  colback=teal!8!white,
  colframe=teal!80!black,
  arc=3mm,
  boxrule=1.1pt,
  width=\linewidth,
  halign=center,
  top=2.2mm, bottom=2.2mm,
  left=3mm, right=3mm,
  drop fuzzy shadow
]
  {\large\bfseries\color{teal!95!black}\faBookmark\ BÀI 71. ÔN TẬP HÌNH HỌC (HÌNH PHẲNG \& HÌNH KHỐI)}\\[2pt]
  {\scriptsize\color{black!75}\itshape Bản đồ công thức Chu vi, Diện tích và Thể tích toàn bộ các hình đã học ở Tiểu học trên 1 trang A4}
\end{tcolorbox}

\vspace{3.2mm}
% --- KHỐI 1: HÌNH PHẲNG VÀ HÌNH KHỐI ---
\noindent
\begin{minipage}[t]{0.52\textwidth}
  \begin{tcolorbox}[
    enhanced,
    title={\footnotesize\bfseries\color{white}\faShapes\ 1. Công thức Hình phẳng},
    colback=blue!2!white,
    colframe=blue!75!black,
    arc=2.5mm, boxrule=0.8pt, top=1.8mm, bottom=1.8mm, left=2.2mm, right=2.2mm
  ]
    \footnotesize
    $\bullet$ \textbf{Tam giác:} $S = (a \times h) : 2$\\
    $\bullet$ \textbf{Hình thang:} $S = (a + b) \times h : 2$\\
    $\bullet$ \textbf{Hình tròn:}\\
    Chu vi: $C = d \times 3,14 = 2 \times r \times 3,14$\\
    Diện tích: $S = r \times r \times 3,14$
  \end{tcolorbox}
\end{minipage}%
\hfill
\begin{minipage}[t]{0.46\textwidth}
  \begin{tcolorbox}[
    enhanced,
    title={\footnotesize\bfseries\color{white}\faCube\ 2. Công thức Hình khối},
    colback=teal!2!white,
    colframe=teal!75!black,
    arc=2.5mm, boxrule=0.8pt, top=1.8mm, bottom=1.8mm, left=2mm, right=2mm
  ]
    \footnotesize
    $\bullet$ \textbf{Hình hộp chữ nhật:}\\
    $S_{xq} = 2(a+b) \times c$; $S_{tp} = S_{xq} + 2ab$\\
    $V = a \times b \times c$\\
    $\bullet$ \textbf{Hình lập phương:}\\
    $S_{xq} = a \times a \times 4$; $S_{tp} = a \times a \times 6$\\
    $V = a \times a \times a$
  \end{tcolorbox}
\end{minipage}

\vspace{3mm}

% --- KHỐI 3: SƠ ĐỒ TƯ DUY ---
\noindent
\begin{tcolorbox}[
  enhanced,
  title={\footnotesize\bfseries\color{white}\faSitemap\ Sơ đồ tư duy: Hệ thống công thức Hình học Tiểu học},
  colback=yellow!2!white,
  colframe=teal!85!black,
  arc=3mm, boxrule=0.9pt, top=1.8mm, bottom=1.8mm, left=1.5mm, right=1.5mm,
  halign=center
]
  \begin{tikzpicture}[
    scale=0.92,
    every node/.style={transform shape},
    root/.style={rectangle, rounded corners=5pt, fill=teal!85!black, text=white, font=\bfseries\footnotesize, inner sep=5pt, align=center},
    b1/.style={rectangle, rounded corners=4pt, fill=blue!12, draw=blue!75!black, text=blue!90!black, font=\bfseries\scriptsize, inner sep=4.5pt, align=center},
    b2/.style={rectangle, rounded corners=4pt, fill=purple!12, draw=purple!75!black, text=purple!90!black, font=\bfseries\scriptsize, inner sep=4.5pt, align=center}
  ]
    \node[root] (core) at (0, 0.85) {\faShapes\ HÌNH HỌC LỚP 5};
    \node[b1] (n1) at (-4.8, -0.45) {\textbf{HÌNH PHẲNG (Chu vi \& Diện tích)}\\[2pt]Tam giác: $S = (a \times h) : 2$\\Hình thang: $S = (a + b) \times h : 2$\\Hình tròn: $C = d \times 3,14$; $S = r^2 \times 3,14$};
    \node[b2] (n2) at (4.8, -0.45) {\textbf{HÌNH KHỐI (Diện tích \& Thể tích)}\\[2pt]HHCN: $S_{xq} = 2(a+b)c$; $V = abc$\\HLP: $S_{xq} = 4a^2$; $S_{tp} = 6a^2$; $V = a^3$\\Hình trụ: $2$ mặt đáy tròn \& $1$ mặt xung quanh};

    \draw[teal!85!black, thick, -{Stealth[length=2mm]}] (core) -| (n1);
    \draw[teal!85!black, thick, -{Stealth[length=2mm]}] (core) -| (n2);
  \end{tikzpicture}
\end{tcolorbox}

\vspace{3mm}

% --- KHỐI 4: QUY TẮC ĐỒNG NHẤT & BẪY TRẮC NGHIỆM ---
\noindent
\begin{minipage}[t]{0.54\textwidth}
  \begin{tcolorbox}[
    enhanced,
    title={\footnotesize\bfseries\color{white}\faCheckDouble\ Quy tắc đồng nhất đơn vị đo},
    colback=green!3!white,
    colframe=green!60!black,
    arc=2.5mm, boxrule=0.8pt, top=1.8mm, bottom=1.8mm, left=2.2mm, right=2.2mm
  ]
    \footnotesize
    $\bullet$ Đáy và chiều cao bắt buộc phải \textbf{cùng một đơn vị đo}.\\
    $\bullet$ Đổi ra cùng đơn vị trước khi tính: $\text{dm} \rightarrow \text{cm}$ hoặc $\text{m}$.\\
    $\bullet$ Đơn vị diện tích có số mũ $2$ ($\text{m}^2, \text{cm}^2$), đơn vị thể tích có số mũ $3$ ($\text{m}^3, \text{dm}^3$).
  \end{tcolorbox}
\end{minipage}%
\hfill
\begin{minipage}[t]{0.44\textwidth}
  \begin{tcolorbox}[
    enhanced,
    title={\footnotesize\bfseries\color{white}\faExclamationTriangle\ Mẹo tránh bẫy trắc nghiệm},
    colback=red!2!white,
    colframe=red!70!black,
    arc=2.5mm, boxrule=0.8pt, top=1.8mm, bottom=1.8mm, left=2.2mm, right=2.2mm
  ]
    \footnotesize
    $\bullet$ Thùng không nắp: $S_{tp}$ chỉ cộng $1$ mặt đáy.\\
    $\bullet$ Nhầm giữa đường kính $d$ và bán kính $r$ ($d = 2 \times r$).\\
    $\bullet$ Bán kính tăng $2$ lần $\rightarrow$ Diện tích tăng $2 \times 2 = 4$ lần!
  \end{tcolorbox}
\end{minipage}

\end{document}
